\section{怎样辨析文言实词的一词多义现象}

中学同学和广大社会青年在学习文言文时，感到古今词义虽然难以辨析，但是同文言文中一词多义现象比较起来，又容易些了，因为现代汉语的词义总是比较熟悉的，关键在于古代汉语中一个文言实词有那么多词义，实在不好掌握。事实的确是这样。

一词多义现象，古今汉语都存在，只是古代汉语更显得突出，而且有些词义与现代的词义相距甚远，不好掌握。古代汉语的一词多义现象之所以显得更加突出，是因为古代文字比现代少得多，古人为了很好地表达自己的思想，只好让一个字承担起多种意义来。但是，这种一词多义现象并不是随心所欲的胡来，而是有规律可循的。这个规律就是一个词总有它本来的意义，这叫作“本义”，而其他的意义是由这个本义引申、生发出来的，这叫作“引申义”。一般说来，一个词不论有多少种意义，用在一个具体的句子中，只能表示这多种意义中的一个。因此，我们在阅读文言文时，一定要根据上下文意来确定某个词的特定词义。

具体说来，我们要辨析文言实词的一词多义现象，可以从以下几个方面着手。

\textbf{1. 掌握中学语文课本中的常见文言实词的词义，是辨析一词多义现象的基础。}

部编中学语文课本的十二册书中，收有古代诗文一百余篇，其中就有常见实词近千个。如果我们能掌握五百个左右的常见实词的各种词义的话，那就有了读懂浅易文言文的基础。而要掌握五百个左右的常见实词的各种词义并不特别困难，因为我们只要认真熟读、细心比较在先秦两汉的名篇中出现的常见实词，也就差不多了。例如：《诗经》的《伐檀》、《硕鼠》，《论语》的《子路、曾皙、冉有、公西华侍坐》，《孟子》的《齐桓晋文之事》，《鱼我所欲也》，《庄暴见孟子》，《庄子》的《庖丁解牛》，《墨子》的《公输》，《荀子》的《劝学》，《左传》的《殽之战》，《曹刿论战》，《吕氏春秋》的《察今》，《战国策》的《邹忌讽齐王纳谏》，贾谊的《过秦论》，《论积贮疏》，司马迁《史记》的《鸿门宴》，《廉颇蔺相如列传》，《信陵君窃符救赵》，《陈涉世家》，以及《韩非子》，《淮南子》，《列子》的寓言故事，此外还有唐宋大散文家韩愈、柳宗元、欧阳修、王安石、苏轼的文章。这些文章中的常见文言实词非常丰富，我们如果能掌握这些常见文言实词的各种词义，就为培养文言文自学能力奠定了基础。

\textbf{2.根据课文中文言实词的一词多义现象，有意识地归纳其本义和引申义。}

我们知道，一个多义词的各种词义之间是有联系的，它们都是由一个基本的词义发展、引申出来的。有的是从本义直接引申出来的，有的是从引申义再进一步引申出来的，前者叫直接引申，后者叫间接引申。我们懂得了这个道理，就能够把一个词的多种意义之间的关系整理得条分缕析，不必要死记一大堆词义而不得要领。例如：
\begin{entry}
	
\item[解]东汉大文学家许慎所著的《说文解字》说：“判也，从刀判牛角。”意思是“剖解牛”。如《庖丁解牛》：“庖丁为文惠君解牛。”这儿用的就是这个词的“本义”。由此可以直接引申为“分解动物的肢体”，如“宰夫将解鼍。”由此再引申为“分解、解体”，如“土崩瓦解”。由“分解动物”的意义还引申为“解绳结”。由“解绳结”的意义再引申为解脱、解说、理解、见解、和解、排解、消除等。我们掌握了“解”字的本义和引申义，就可以根据上下文意来确定它在文句中的特定含义。

\item[兵]许慎《说文解字》曰：“械也，从甘持斤。”意思是“兵器”，字形象是用双手拿着斧斤之类的兵器。如《六国论》：“六国破灭，非兵不利，战不善，弊在赂秦。”《过秦论》：“斩木为兵，揭竿为旗。”都是用的这个字的本义。由这个本义可以引申为“持兵的人”，即“士兵”，如《鸿门宴》：“沛公兵十万，在霸上。”又可以引申为“军队”，如《信陵君窃符救赵》：“公子遂将晋鄙军，勒兵，下令军中。”由此又可以引申为“战阵之事、军事、战争”，如《史记·孙子吴起列传》：“兵法：百里而趣利者蹶上将······”又说：“庞涓自知智穷兵败，乃自刭。”

\item[引]许慎《说文解字》说：“引，开弓也。”从字形构造上看，从“弓”从“丨”（“丨”就是箭），意思是说，把箭搭上弦，慢慢把弓向后拉开，这就叫“引”。如毛主席在《湖南农民运动考察报告》中引用《孟子》的话：“引而不发，跃如也。”（拉开了弓不射出去，却摆着跃跃欲动的姿势。）用的就是这“引”字的本义。由这个本义直接引申为“拉开、拉出”，如《廉颇蔺相如列传》：“左右或欲引相如去。”又因拉开弓而使弓和弦的距离伸长，所以引申为“伸长”、“延长”，如《中山狼传》：“狼奄至，引（伸长）首顾曰······”《三峡》：“常有高猿长啸，属引（延长）凄异。”由于开弓而使弓弦向后退，所以又可以引申为“向后退”、“避开”，如《中山狼传》：“引避道左，以待赵人之过。”由于开弓而指向目标，所以又引申为“引导”，再引申为“带领”，如《信陵君窃符救赵》：“平原君负魏失为公子先引（引导）。”《赤壁之战》：“操乃······引（带领）军北还。”

\item[亡]许慎《说文解字》说：“逃也，从人从L。”“L”即古“隐”字。从人从L，意思说人在回曲隐蔽的地方。这也就是说，“亡”的本义是“逃跑”，如《陈涉世家》：“今亡亦死，举大计亦死，等死，死国可乎？”特指“出奔”，“逃到国外去”，如《廉颇蔺相如列传》：“臣尝有罪，窃计欲亡走燕。”引申为“失掉”（让它跑掉），如《屈原列传》：“兵挫地削，亡其六郡，身客死于秦，为天下笑。”由失掉、消失而引申为“死”（与“存”相对），如范缜《神灭论》：“岂容形亡而神在？”再引申为“灭亡”，如《过秦论》：“山东豪俊遂并起而亡秦族矣。”

\item[逆]许慎《说文解字》说：“迎也，从辵声。关东曰逆，关西曰迎。”“辵”即古“踏”字，“乍行乍止”的意思；“逆”读yì。从字形构造上看，“逆”是个形声字，“迎”是“逆”的本义。如《送东阳马生序》：“寓逆旅，主人日再食。”“逆旅”指迎接旅客的房舍，这儿用的正是本义。再如《赤壁之战》：“将兵与备并力逆操。”“逆”，迎击之义。由“迎”引申为“预先”（限用于“预见”的意义上），如《后出师表》：“凡事如此，难可逆见。”再引申为“向相反方向活动”，“倒着”，如《水经注·江水》：“水逆流百余里，涌起数十丈。”再引申为“不顺”，“抵触”（与“顺”相对），如《史记·留侯世家》：“忠言逆耳利于行。”由此再引申为“背叛”，“叛逆”，如《指南录后序》：“予自度不得脱，则直前诟虏帅失信，数吕师孟叔侄为逆。”
\end{entry}

以上数例说明：要辨析文言单词的本义，主要靠字形构造的分析，并寻找古代文献的用例。因为汉字是表意体系的文字，一个字的构造往往就是理解这个字的钥匙。所以，要初学古代汉语的人从字形构造上辨析一个字的本义是相当困难的。一是由于这些汉字的字形笔势已经与甲骨文、金文等古文字字形相比，起了很大的变化，难以看出原来的字形结构，二是由于中学生和一般青年同志还没有接触过文字学，即使把一个字的原形摆在面前，也无法辨析。我们只有一个可行的办法，那就是为了弄清一个文言单词的本义和引申义，在查阅《辞源》、《辞海》、《康熙字典》等工具书时，认准一个字的第一个义项（对词义的第一个解释）一般就是本义，然后再从本义仔细体会引申义。这虽是一个笨办法，对我们来说却是切实有效的。

至于一个文言实词除了本义和引申义外，还有假借义，这留到下一节中去谈。

\textbf{3. 根据上下文意，确定一个文言多义词在句中的特定意义。}

一词多义只是就一个个孤立的文言实词而言，事实上用在句子中的文言实词只能有一个特定的意义。

怎样确定一个文言多义词在句中的特定意义呢？一是要掌握这个多义词的各个词义，包括本义和引申义，有的还有假借义；二是要根据“词不离句”的原则，辨析这个多义词在句中的词性和含义。例如：

宋有澄子者，亡缁衣，求之涂（途）。见妇人衣缁衣，援而弗舍，欲取其衣。曰：“今者我亡缁衣。”妇人曰：“公虽亡缁衣，此实吾所自为也。”澄子曰：“子不如速与我衣。昔我所亡者，纺缁也；今子之衣，禅缁也。以禅缁当纺缁，子岂不得哉？”（《吕氏春秋·淫辞》）

[\ 缁（zī）衣：帛制的黑色衣服。纺缁：有里子的缁衣。禅缁：没有里子的缁衣。\ ]

这段短文中加点的“亡”、“衣”、“援”、“舍”、\\
“公”、“为”、“子”等，都是常见文言多义词。“亡”的义项很多，但从上下文意看，这儿只能作“丢失”讲。“衣”的本义是“上衣”（与“裳”相对），“缟衣”的“衣”就是用的本义，引申为“衣服”，都是名词；而文中加点的“衣”字却是动词，处在谓语位置上，因此就要用另一个义项“穿（衣服）”来解释才行。“援”的本义是“引”（开弓），引申为“拉、拽”，再引申为“领来”，再引申为“引用、引证”，再引申为“帮助、救助”，从上下文意来判断，这儿只能作“拉、拽”讲。“舍”的本义是“宾馆”，引申为“房舍”，是名词；由此引申为动词“住一夜”，特指行军或狩猎的临时住宿；又一个义项是军行三十里为一舍，另外一个义项是动词“放弃、不要、不敢”（这种意义后来写作“擒”，现在又简化成“舍”），由此引申为“释放”、“放开”。从这段短文的上下文意看，这儿应作“放开”讲。“公”的本义是“公家的”，与“私人的”相对；又作名词“公家”讲，引申为“公正、无私”；又作副词“公然、公开地”讲；又作公、侯、伯、子、男“五等爵位的第一等”；又作“官职的最高级，在卿之上”讲；由此引申为对人的敬称，相当现代的“您”，这段文字中的“公”正用这个引申义。“为”字义项很多，但根据本句上下文意可以判断这个“为”字是动词，动词“为”的基本意义是“做、造作”，引申为“当作、作为”，又可引申为“变成、成为”，还可引申为“算做、算是”，这段文字中正是用的动词“为”的基本意义“做”。“子”的本义是“婴儿”，引申为“儿女”，一般指儿子，有时指女儿，又引申为“植物的籽实”，它的又一个义项是“男子的尊称”，用来专指有德之人，相当于“夫子”；由此引申为对人的尊称，相当于现代汉语的“您”，这段文字中的“子”正是用这个引申义“您”，它不仅尊称男人，也尊称女人。

除了上面辨析的几个常见文言多义词之外，我们再列举六十个常见文言多义词，说明它们的主要义项，作为本节的“附录”。请读者同志从中学语文课本中选取例句来印证这些主要义项，并请您照此办法整理和积累常见文言多义词的主要义项。

总之，我们只要熟练掌握课本中文言实词的主要义项，又能有意识地归纳它们的本义和引申义，再在具体文言短文中根据“词不离句”的原则分辨其特定的意义，那么经过一段时间的积累和练习，自能培养好自学文言文的能力。
\newpage